\documentclass[10pt,a4paper]{amsart}
\usepackage[margin=19mm]{geometry}
\usepackage{amsmath,amssymb,mathtools}
\usepackage[scr=rsfs,cal=euler]{mathalfa}
\usepackage{xfrac}
\usepackage[unicode,hidelinks,bookmarks=false]{hyperref}
\newcommand{\ie}{i.e.\ }
\newcommand{\cf}{cf.\ }
\newcommand{\resp}{respectively\ }
\newcommand{\Subsection}{\subsection}
\providecommand{\nobreakdash}{\nobreak\hbox{-}}
\setlength{\emergencystretch}{2em}
\multlinegap0pt
\usepackage{bm}
\usepackage{bbm}
\usepackage{dsfont}
\usepackage{xspace}
\usepackage{cancel}
\usepackage{mathtools}
\usepackage{amsthm}
\makeatletter
\@ifundefined{theorem}{\newtheorem{theorem}{Theorem}}{}
\makeatother
\makeatletter
\expandafter\ifx\csname Assumptions\endcsname\relax
\newenvironment{Assumptions}% Definition of assumptions
{%
\setcounter{enumi}{0}
\renewcommand{\theenumi}{(\textbf{A}.\arabic{enumi})}
\renewcommand{\labelenumi}{\theenumi}
\begin{enumerate}}%
\fi
\newcommand{\R}{\mathbb{R}}
\newcommand{\g}[1]{\mathcal{G}_{#1}}
\newcommand{\gk}[1]{\mathbb{G}_{#1}}
\makeatother
\title[Large Solutions for Fractional Laplacian on Infinite Cylindr]{Large Solutions for Fractional Laplacian on Infinite Cylindrical Domains}
\author{Indranil Chowdhury, N N Dattatreya}
\date{Source: arXiv:2607.25906v1 (CC BY 4.0)}
\begin{document}
\maketitle

\noindent\emph{Source and licence.} Large Solutions for Fractional Laplacian on Infinite Cylindrical Domains. Version arXiv:2607.25906v1 (CC BY 4.0), distributed under the Creative Commons Attribution 4.0 International licence, as recorded in the arXiv licence metadata for this exact version and shown on the abstract page https://arxiv.org/abs/2607.25906v1. Full rights notice: licenses/arxiv-2607.25906-CC-BY-4.0.txt.\par\smallskip

\noindent\textit{Editorial scope.} Counted unit: the Theorem whose source label is \texttt{theorem greens poteantial continuity on unbounded domains} in the article (source0.tex lines 658--669, source label \texttt{theorem greens poteantial continuity on unbounded domains}), delivered together with the article's own complete original proof of it (source0.tex lines 670--722, one \texttt{proof} environment of 53 lines). No mathematics is written, repaired or re-derived: statement and proof are the article's own text line for line. Required context retained uncounted: eqn greens estimatte on unbounded domain (lines 520--522); eqn behaviour of greens on unbounded domains (lines 535--537). Every definition, display and label of those lines is retained unchanged. Omitted from this package and not redistributed: 1--519, 523--534, 538--657, 723--2059. Nothing inside a retained excerpt is omitted or altered: every retained body is delivered exactly as the article's lines are written, so no omission receipt of any kind accompanies this package. Citations: the retained excerpts carry no citation command at all, so no bibliography entry of the article is transcribed and no reference is redistributed. Reference adaptation: none was needed and none was made; every reference inside the delivered unit resolves inside it, so no unresolved reference or citation is produced. Printed slips kept literal: no printed word, symbol, hypothesis or number of the counted statement or of its proof was altered. Layout and class: the article's own class is \texttt{amsart} with packages including graphicx, inputenc, amsmath, xcolor, amssymb, stackengine, amsfonts, amstext, amsthm, babel, csquotes, nth, hyperref, biblatex; the wrapper is the lane's pinned \texttt{amsart} editorial wrapper at 10pt on A4, which restates the theorem-like environments the retained text needs and adds the article's own macro definitions that the retained text needs (\texttt{\textbackslash R}, \texttt{\textbackslash g}, \texttt{\textbackslash gk}), with extra wrapper packages (bm, bbm, dsfont, xspace, cancel, mathtools). The delivered numbering is the unit's own. Assets: no retained span contains a figure environment, a graphics-inclusion command, a PDF-page inclusion, a table or a TikZ picture; no image file is copied into this package, the package declares no asset dependency and no image file is redistributed with it. Licence: version arXiv:2607.25906v1 of this article is distributed under the Creative Commons Attribution 4.0 International licence, as recorded in the arXiv licence metadata for this exact version and shown on the abstract page https://arxiv.org/abs/2607.25906v1. A full rights notice is delivered at licenses/arxiv-2607.25906-CC-BY-4.0.txt. Characters: every retained excerpt is pure ASCII, so the delivered engine, XeLaTeX with its default Latin Modern text font, reports no missing character.
\begin{equation}\label{eqn greens estimatte on unbounded domain}
    \gk{\Omega}(x,y)\leq C(n,s)|y-a_x|^{s}\frac{\delta_{\Omega}^{s}(x)}{|x-y|^{n-s}}.
\end{equation}

   \begin{equation}\label{eqn behaviour of greens on unbounded domains}
           \gk{\Omega}(x,y)\leq \gk{B^{C}}(x,y) \leq c|y-a_x|^{s}\frac{\delta_{B}^{s}(x)}{|x-y|^{n-s}}=c|y-a_x|^{s}\frac{\delta_{\Omega}^{s}(x)}{|x-y|^{n-s}}. \qedhere
   \end{equation}

\begin{theorem}\label{theorem greens poteantial continuity on unbounded domains}
    Let $\Omega$ be a domain in $\R^{n}$ such that $|\R^{n}\setminus \Omega|>0$, and has $C^{1,1}$ boundary. The map $\g{\Omega}$ is continuous between
    \begin{equation*}
    \begin{split}
        L^{\infty}_{c}(\Omega)\to \delta_{\Omega}^{s}L^{\infty}(\Omega),\quad \text{and} \quad L^{1}(\Omega,\delta_{\Omega}^{s})\to L_{loc}^{1}(\Omega). 
    \end{split}    
    \end{equation*}
Further, for $q>n/2s$, it is a continuous map between 
\begin{equation*}
    L^{1}(\Omega,\delta_{\Omega}^{s})\cap L_{loc}^{q}(\Omega)\to L^{\infty}_{loc}(\Omega).
\end{equation*}
\end{theorem}

\begin{proof}
\textit{First part:}
Let $f\in L^{\infty}_{c}(\Omega)$. Denote $K=$ Supp$(f)$. Choose $\rho>0$ such that $K\subset\subset\Omega\setminus\Omega_{\rho}=\{x\in \Omega~|~\delta_{\Omega}(x)\geq\rho\}$. Then for any $x\in \Omega\setminus \Omega_{\rho}$ %such that $K \subset K_{1}\subset\subset\Omega$, since $\delta_{\Omega}(x)\geq d(K_{1},\partial\Omega)$, we obtain by 
using $\gk{\Omega}(x,y)\leq \Gamma(x-y)$ that
\begin{equation*}
        \g{\Omega}(f)(x)\leq \|f\|_{L^{\infty}(\Omega)}\int_{K}\gk{\Omega}(x,y)\ dy
        \leq  c(n,s)\|f\|_{L^{\infty}(\Omega)}\int_{K}\frac{dy}{|x-y|^{n-2s}}\leq C(K,\rho,n,s)\|f\|_{L^{\infty}(\Omega)} \delta_{\Omega}^{s}(x).
    \end{equation*}
    Indeed, if $B_{1}(x)\cap K\neq \emptyset$ then the integral over $K$ is bounded by integral on a ball of radius $1+$diam$(K)$. Otherwise the integral is bounded by $|K|$.

    
    For $x\in \Omega_{\rho}$, using \eqref{eqn behaviour of greens on unbounded domains}
     \begin{equation}\label{eqn greens potential for supported f}
        \g{\Omega}(f)(x)\leq \|f\|_{L^{\infty}(\Omega)}\int_{K}\gk{\Omega}(x,y)\ dy
        \leq  c(n,s)\|f\|_{L^{\infty}(\Omega)}\int_{K}|y-a_x|^{s}\frac{\delta_{\Omega}^{s}(x)}{|x-y|^{n-s}}\ dy.
    \end{equation}
Further, by triangle inequality
    \begin{equation}\label{eqn triangle inequality}
        \frac{|y-a_x|}{|x-y|}\leq 1+\frac{|x-a_x|}{|x-y|}\leq 1+\frac{|x-a_x|}{d(K,\Omega_{\rho})}.
    \end{equation}
     Due to the uniform ball condition, $|x-a_x|\leq \rho+R$. Using these in \eqref{eqn greens potential for supported f}
    \begin{equation*}
       \g{\Omega}(f)(x)\leq C(K,\rho,n,s)\|f\|_{L^{\infty}(\Omega)}\delta_{\Omega}^{s}(x)\int_{K}\frac{dy}{|x-y|^{n-2s}}\leq C(K,\rho,n,s)\|f\|_{L^{\infty}(\Omega)} \delta_{\Omega}^{s}(x).
    \end{equation*}
    
    \textit{Second part:}
    Let $f\in L^{1}_{loc}(\Omega,\delta_{\Omega}^{s})$ and let $K$ be any compact subset of $\Omega$. Choose $\rho>0$ such that $K\subset\subset \Omega\setminus\Omega_{\rho}$. Then by using $\gk{\Omega}(x,y)\leq \Gamma(x-y)$ and \eqref{eqn greens estimatte on unbounded domain}, we obtain
    \begin{equation*}
        \begin{split}
            \int_{K}|\g{\Omega}(f)(x)|\ dx&=\int_{K}\left\{\int_{\Omega\setminus\Omega_{\rho}}+\int_{\Omega_{\rho}}\right\}|f(y)|\gk{\Omega}(x,y)\ dy\ dx\\
            &\leq \int_{\Omega\setminus\Omega_{\rho}}\int_{K}\frac{|f(y)|}{|x-y|^{n-2s}}\ dx\ dy+\int_{\Omega_{\rho}}|f(y)|\delta_{\Omega}^{s}(y)\int_{K}\frac{|x-a_y|^{s}}{|x-y|^{n-s}}\ dx \ dy.
        \end{split}
    \end{equation*}
    By using the triangle inequality as in \eqref{eqn triangle inequality}, for the second inequality, we infer that
    \begin{equation*}
        \int_{K}|\g{\Omega}(f)(x)|\ dx\leq C(K,n,s,\rho)\|f\|_{L^{1}(\Omega,\delta_{\Omega}^{s})}+\Tilde{C}(K)\|f\|_{L^{1}(\Omega,\delta_{\Omega}^{s})}\leq C(K,n,s)\|f\|_{L^{1}(\Omega,\delta_{\Omega}^{s})}.
    \end{equation*}

    \textit{Third part:} Let $f\in L^{1}(\Omega,\delta_{\Omega}^{s})\cap L_{loc}^{q}(\Omega)$, for some $q>n/2s$, and $K$ be a compact subset of $\Omega$. For $\rho>0$ as in the previous case, then, as in the last part
    \begin{equation*}
        |\g{\Omega}(f)(x)| \leq \int_{\Omega\setminus\Omega_{\rho}}\frac{|f(y)|}{|x-y|^{n-2s}}\ dy+\int_{\Omega_{\rho}}|f(y)|\delta_{\Omega}^{s}(y)\frac{|x-a_y|^{s}}{|x-y|^{n-s}} \ dy,
    \end{equation*}
    Again, by the triangle inequality as in \eqref{eqn triangle inequality} %$\frac{|x-a_y|}{|x-y|}\leq 1+\frac{|y-a_y|}{r}$
    , we have
    \begin{equation*}
        \int_{\Omega_{\rho}}|f(y)|\delta_{\Omega}^{s}(y)\frac{|x-a_y|^{s}}{|x-y|^{n-s}} \ dy\leq C(K,\rho)\|f\|_{L^{1}(\Omega,\delta^{s}_{\Omega})}.
    \end{equation*}
    Choose $r>0$ such that for all $x\in K$, $B_{r}(x)\subset\Omega$, by H\"older inequality
    \begin{equation*}
        \int_{(\Omega\setminus\Omega_{\rho})\cap B_{r}(x)}\frac{|f(y)|}{|x-y|^{n-2s}}\ dy\leq \|f\|_{L^{q}(B_{r}(x))}\left\{\int_{B_{r}(x)}\frac{dy}{|x-y|^{q'(n-2s)}}\right\}^{1/q'},
    \end{equation*}
    and is finite owing to $q>n/2s$. The result follows by combining the above equations with the integral on $(\Omega\setminus\Omega_{\rho})\setminus B_{r}(x)$, which is bounded by $\left(1/(r^{n-2s}\rho^{s})\right)\|f\|_{L^{1}(\Omega,\delta_{\Omega}^{s})}$.
\end{proof}

\end{document}
